% methods.tex (v12, 2026-10-18: Fig. 2 caption follows the redesigned figure): Data and methods, placed before Results. Unnumbered subsections; equations (1)-(9).
\section{Data and methods}\label{sec:methods}

\subsection{Study regions and active-layer observations}
We used three regions with dense, multi-year thaw-depth observations: Alaska, the Lena Delta and north-western Canada (Fig.~\ref{fig:data}, Table~\ref{tab:data}). Observations for Alaska and Canada were taken from the ABoVE soil moisture and active-layer thickness compilation, version 2\cite{Moore2025}, and from the CALM end-of-season grid values\cite{Streletskiy2025}; observations for the Lena Delta were taken from the ALLena thaw-depth data set\cite{Veremeeva2025}. ABoVE probe measurements were restricted to August--October and to $0<\mathrm{ALT}<300$ cm, because measurements made earlier in the season record thaw depth rather than the seasonal maximum. ALLena measurements were restricted to August--September values without the probe-limit flag, and the seasonal maximum was taken for each location and year. All observations were assigned to 1 km cells and averaged within each cell and year, retaining the mean day of observation; CALM values without a date were assigned day 244 (1 September). We kept the years 2010--2024, for which year-specific climate inputs are available. The resulting labels comprise \nncyAk{} cell-years in Alaska (\nyrsAk), \nncyLe{} in the Lena Delta (\nyrsLe) and \nncyCa{} in Canada (\nyrsCa), distributed over \nnblkAk, \nnblkLe{} and \nnblkCa{} blocks of 0.5$^\circ$ (Table~\ref{tab:data}). Sources, versions and licences are listed in Supplementary Table~S1, the years used in each analysis in Supplementary Table~S2 and the processing steps in Supplementary Methods~S1.

% table1.tex (v11): Table 1, labels by region and the periods of inputs and products. Numbers from macros.
\begin{table}[!htbp]
\caption{Data and years used. Top: ALT labels by region (cell-years of 2010--2024 with year-specific climate input). Bottom: inputs and products, with the period available and the years used in this study. Blocks are the 0.5$^\circ$ cells used for spatial cross-validation and bootstrap; there are no Canadian labels in 2020.}\label{tab:data}
\scriptsize
\begin{tabular}{@{}lrrrlrl@{}}
\toprule
Region & Cell-years & Cells & Blocks & Years & Cells $\geq$4 yr & Largest year \\
\midrule
Alaska & \nncyAk{} & \nncellAk{} & \nnblkAk{} & 2010--2024 & \nnmultiAk{} & \ntopyrAk{} (\ntopyrshareAk\%) \\
Lena Delta & \nncyLe{} & \nncellLe{} & \nnblkLe{} & 2010--2022 & \nnmultiLe{} & \ntopyrLe{} (\ntopyrshareLe\%) \\
Canada (NW) & \nncyCa{} & \nncellCa{} & \nnblkCa{} & 2010--2024 & \nnmultiCa{} & \ntopyrCa{} (\ntopyrshareCa\%) \\
External CALM & \nftScoredRows{} & \nftScoredCells{} & \nftScoredBlocks{} & \nyrsext{} & -- & -- \\
\bottomrule
\end{tabular}

\vspace{6pt}
\begin{tabular}{@{}>{\raggedright\arraybackslash}p{0.25\textwidth}>{\raggedright\arraybackslash}p{0.12\textwidth}>{\raggedright\arraybackslash}p{0.13\textwidth}>{\raggedright\arraybackslash}p{0.16\textwidth}>{\raggedright\arraybackslash}p{0.16\textwidth}@{}}
\toprule
Data & Resolution & Available & Used & Role \\
\midrule
ERA5-Land monthly\cite{munozsabater2021_era5_land} & 0.1$^\circ$ & 1950-- & 2010--2024; mean 2011--2024 & Climate input \\
Copernicus DEM\cite{CopDEM2021} & 30 m & static & -- & Terrain \\
SoilGrids 2.0\cite{Poggio2021} & 250 m & static & -- & Soil \\
MODIS MOD13Q1, MOD10A2\cite{Didan2021,HallRiggs2021} & 250--500 m & 2000-- & 2010--2024 & Sensitivity test \\
ESA CCI permafrost probability\cite{Westermann2024} & 1 km & 1997--2021 & mean & Map domain \\
Wei et al.\cite{Wei2026} & 1 km & 2000--2024 & \nyrsprodWei & Comparison \\
Liu et al.\cite{Liu2024} & 1 km & 2003--2020 & \nyrsprodLiu & Comparison \\
ESA CCI ALT v4\cite{Westermann2024} & 1 km & 1997--2021 & \nyrsprodCcifour & Comparison \\
ESA CCI ALT v5\cite{Westermann2025} & 1 km & 1997--2023 & \nyrsprodCcifive & Comparison \\
Yi and Kimball\cite{YiKimball2019} & 1 km & 2001--2015 & \nyrsprodYk{} (Alaska) & Comparison \\
Aalto et al.\cite{Aalto2018} & 30$''$ & 2000--2014 & time-invariant & Comparison \\
\bottomrule
\end{tabular}
\end{table}


\paperfig{Fig1_data}{1.0}{\textbf{Study regions, observations and data periods.} \textbf{a}, ESA CCI permafrost probability (mean of 1997--2021) with the three study regions and the number of 1 km cells with ALT labels in 2010--2024. \textbf{b}, Number of cell-year labels per observation year, 2010--2024. \textbf{c}, Periods covered by the labels, the ERA5-Land monthly data used for the climate input, the static covariates and the six published products; for each product the period available and the years used for comparison at the label cells are distinguished, and the shaded band marks the analysis period 2010--2024.}{fig:data}

\subsection{Year-specific climate input}
Monthly ERA5-Land fields for January 2010 to December 2024\cite{munozsabater2021_era5_land} (2 m air temperature, top-layer soil temperature and snow water equivalent) were extracted at the nearest land grid cell of each 1 km cell. For cell $c$ and year $y$, the thawing degree-day sum is
\begin{equation}
\mathrm{TDD}_{c,y} = \sum_{m=1}^{12} \max\!\left(T_{c,y,m},\,0\right) d_{y,m},
\label{eq:tdd}
\end{equation}
where $T_{c,y,m}$ is the monthly mean air temperature and $d_{y,m}$ the number of days in month $m$. Because most observations were made before the end of the thaw season, the model input is the TDD accumulated to the day of observation $t$,
\begin{equation}
\mathrm{TDD}^{\mathrm{obs}}_{c,y} = C_{c,y,k-1} + \frac{t - t_{k-1}}{d_{y,k}} \left( C_{c,y,k} - C_{c,y,k-1} \right),
\label{eq:tddobs}
\end{equation}
where $C_{c,y,k}$ is the cumulative TDD to the end of month $k$, $k$ is the month containing $t$ and $t_{k-1}$ is the last day of the preceding month. Seven further climate variables were computed for each cell and year: the mean annual air temperature, $\sqrt{\mathrm{TDD}}$, the warmest and the coldest monthly mean temperature, the freezing degree-days of the preceding winter (October $y-1$ to April $y$), the mean top-layer soil temperature and the mean snow water equivalent of the hydrological year (October $y-1$ to September $y$) (Supplementary Table~S3, Supplementary Methods~S2). Two alternative inputs were evaluated on the same cell-years: the full-season TDD of the observation year, and the 2011--2024 mean of the yearly values (mean climate input), which corresponds to the multi-year inputs of existing products. Yearly maps use the full-season TDD of each year and therefore represent end-of-season ALT.

\subsection{Static covariates}
Six terrain variables (elevation, slope, the sine and cosine of aspect, topographic position index and roughness) were computed from the 30 m Copernicus DEM\cite{CopDEM2021} in a $33\times33$ pixel window (about 1 km) around each cell. Nine soil variables (clay, sand, silt, bulk density, coarse fragments and pH at 5--15 cm, and soil organic carbon at 0--5, 5--15 and 15--30 cm) were sampled from SoilGrids 2.0 at 250 m\cite{Poggio2021}. The main feature vector $\mathbf{x}$ therefore contains 23 variables (eight climate, six terrain and nine soil). Sensitivity tests used nested sets: climate only; climate and terrain; the main set plus MODIS summer NDVI, EVI and snow-off timing for 2010--2024 (MOD13Q1 and MOD10A2, version 6.1)\cite{Didan2021,HallRiggs2021}; the main set plus high-resolution soil carbon and land-cover classes\cite{Palmtag2022}; and all variables including the climate of the preceding year. The ESA CCI permafrost probability (mean of 1997--2021)\cite{Westermann2024} was used only to define the map domain (probability of at least 50\%) and not as a predictor.

\subsection{Model formulation}
The Stefan model describes the thaw depth as
\begin{equation}
\widehat{\mathrm{ALT}}_{\mathrm{S}} = E\,\sqrt{\mathrm{TDD}^{\mathrm{obs}}},
\label{eq:stefan}
\end{equation}
where $E$ (cm\,($^\circ$C\,d)$^{-1/2}$) summarises the thermal properties and moisture of the soil\cite{Stefan1891,Riseborough2008}. The relation holds within a season at the Lena Delta monitoring sites and, with more scatter, among years at cells observed repeatedly (Supplementary Fig.~S1). For a set of $n$ training labels $y_i$ with $s_i = \sqrt{\mathrm{TDD}^{\mathrm{obs}}_i}$, the coefficient is recalibrated by least squares through the origin,
\begin{equation}
\hat{E} = \frac{\sum_{i=1}^{n} y_i\, s_i}{\sum_{i=1}^{n} s_i^{2}}.
\label{eq:E}
\end{equation}
The source-coefficient Stefan model uses $\hat{E}$ fitted on the other two regions, and the recalibrated Stefan model uses $\hat{E}$ fitted on the training labels of the target region. The physics-residual model retains the recalibrated Stefan prediction as an anchor and adds a learned correction (Fig.~\ref{fig:workflow}):
\begin{equation}
\widehat{\mathrm{ALT}} = \hat{E}\, s + \lambda\, g(\mathbf{x}), \qquad g = \arg\min_{f} \sum_{i=1}^{n} \left[ \left( y_i - \hat{E}\, s_i \right) - f(\mathbf{x}_i) \right]^{2},
\label{eq:residual}
\end{equation}
where $g$ is a CatBoost gradient-boosting regressor\cite{Prokhorenkova2018} (600 iterations, learning rate 0.03, depth 6, L2 regularisation 3) and $\lambda=0.5$ shrinks the correction towards the physical anchor ($\lambda = 0.25$, 1 and a value selected by inner block cross-validation were tested as sensitivity cases). Three alternative structures place the physical information elsewhere: direct ML fits CatBoost to ALT from $\mathbf{x}$ alone; physics-input ML adds the recalibrated Stefan and Kudryavtsev\cite{Kudryavtsev1974} predictions as features; and physics pseudo-labels augment the training set with an equal number of rows drawn from the cells of the scoring fold, with their covariates and the training-fold Stefan prediction as the target. All learned models were fitted with three random seeds and their predictions averaged. Random forest, a multilayer perceptron, TabM and FT-Transformer\cite{Gorishniy2021,Gorishniy2025,Holzmuller2024} replaced CatBoost in a learner sensitivity test, and SHAP values\cite{Lundberg2017} of the residual and direct models were computed for models fitted to all labels of each region.

\paperfig{Fig2_workflow}{1.0}{\textbf{Cell-year design, Stefan anchor, physics-residual model and validation designs.} \textbf{a}, Each label is one 1 km cell in one year. Left, thaw depth on the observation day and ALT at the end of the thaw season. Centre, cumulative thawing degree-days of a Lena Delta cell in 2013 (black), the 2011--2024 mean of the cell (dashed) and its other years (grey); the model input is the sum accumulated to the observation day (TDD$_{\mathrm{obs}}$, open circle) and the yearly maps use the full-season sum (filled circle). Right, observed ALT against $\sqrt{\mathrm{TDD_{obs}}}$ for the \nncyAk{} Alaska cell-years, with the Stefan anchor fitted by least squares through the origin (equation~(\ref{eq:E})) and five residuals highlighted. \textbf{b}, Physics-residual model: the Stefan anchor and gradient boosting of its residual from 23 climate, terrain and soil inputs, added with the weight $\lambda = 0.5$ (equation~(\ref{eq:residual})); the maps use the full-season TDD of 2014 (Alaska). \textbf{c}, Validation designs on the Lena Delta label cells (one fold each; test year 2016), and ALT anomalies of the Alaska map in 2012 and 2019 relative to the 2011--2024 mean.}{fig:workflow}

\subsection{Validation designs and uncertainty}
All models were evaluated on out-of-fold predictions under five designs\cite{Roberts2017,Ploton2020}: random five-fold cross-validation by cell-year; spatial-block five-fold cross-validation, with 0.5$^\circ$ blocks assigned to folds balanced by the number of cell-years (main design); year hold-out, in which each observation year with at least 20 cell-years is predicted by models trained on the other years; block-by-year hold-out, in which both the block fold and the year of the scored cell-years are excluded from training; and region hold-out, in which a region is predicted by models trained on the other two regions after removing training cells within 100 km of the test region. Accuracy is expressed as
\begin{equation}
\mathrm{RMSE} = \left[ \frac{1}{n} \sum_{i=1}^{n} \left( \hat{y}_i - y_i \right)^{2} \right]^{1/2}, \qquad \Delta\mathrm{RMSE}_{A,B} = \mathrm{RMSE}_{A} - \mathrm{RMSE}_{B},
\label{eq:rmse}
\end{equation}
where both terms of $\Delta\mathrm{RMSE}$ are computed on the same cell-years, so that negative values favour method $A$. Uncertainty was quantified with a paired block bootstrap: the 0.5$^\circ$ blocks were resampled with replacement $B=2000$ times (1000 times for the label-number experiments), $\Delta\mathrm{RMSE}$ was recomputed from all cell-years of the resampled blocks, and the 2.5th and 97.5th percentiles of the replicates define the 95\% interval,
\begin{equation}
\mathrm{CI}_{95} = \left[ Q_{0.025}\!\left( \Delta\mathrm{RMSE}^{(b)} \right),\; Q_{0.975}\!\left( \Delta\mathrm{RMSE}^{(b)} \right) \right], \quad b = 1,\ldots,B.
\label{eq:ci}
\end{equation}
We describe a difference as resolved when its interval excludes zero. Conditions of the gain were examined within regional terciles of eight covariates (full-season TDD of the observation year, freezing degree-days, snow water equivalent, soil organic carbon, clay, slope, elevation and summer NDVI); terciles with fewer than 30 cell-years or fewer than five blocks were not evaluated. Year-wise RMSE was computed for years with at least 30 cell-years, and the correlation between observed and predicted year-to-year changes was computed at cells observed in several years.

\subsection{Label-number and transfer experiments}
Within each region, the training labels of each spatial-block fold were subsampled to $n$ = 25, 50, 100, 250, 500, 1000, 2500 or all (three random draws), the Stefan coefficient was recalibrated on the same $n$ labels, and the complete test fold was scored. For transfer, each region was scored in its spatial-block folds by models trained on all labels of the other two regions plus $n$ = 0, 25, 50, 100, 250 or 500 labels drawn from the training folds of the target region; $\hat{E}$ was refitted on the $n$ target labels (for $n\geq3$), and the residual learner was trained on the source rows, anchored with the source coefficient, together with the target rows, anchored with the refitted coefficient. The external test used CALM observations of 2010--2024\cite{Streletskiy2025} outside the three study regions and outside mountain permafrost (Switzerland and Kazakhstan). Censored values, values flagged as early, burnt, flooded or not reaching the permafrost table, and sites within 100 km of the study labels were removed, leaving \nextcy{} cell-years in \nextcells{} cells in eight regions (Supplementary Methods~S3). Each external region was scored block by block (cell by cell for Greenland, which forms one block) by models trained on the three study regions plus $n$ = 0, 5, 10, 25, 50 or all labels from the other blocks of the same region. One Svalbard borehole with a mean ALT of about 6 m, outside the training range, was excluded from scoring, which left \nftScoredRows{} cell-years in \nftScoredCells{} cells and \nftScoredBlocks{} blocks. The runs are listed in Supplementary Table~S4.

\subsection{Comparison with published products}
Six products were sampled at the label cells: Wei et al.\cite{Wei2026} (1 km, 2000--2024), Liu et al.\cite{Liu2024} (1 km, 2003--2020), ESA CCI Permafrost ALT v4 (1997--2021) and v5 (1997--2023)\cite{Westermann2024,Westermann2025}, Yi and Kimball\cite{YiKimball2019} (Alaska, 2001--2015) and Aalto et al.\cite{Aalto2018} (one value per cell for the 2000--2014 baseline). For yearly products the value of the observation year was used (Supplementary Methods~S4). Products and models were compared on common cell-years, at which Wei et al., Liu et al. and ESA CCI v4 all provide a value for the observation year (\nyrscc). Wei et al. published the coordinates of their training sites, and a leakage mask removed label cells within 25 km of these sites. Liu et al. were trained on an earlier version of the ABoVE observations used here\cite{Liu2024}, so their Alaska and Canada values are not independent of our labels. To test whether the residual learner can also correct a product, the product value $p$ of the observation year was recalibrated on the training fold,
\begin{equation}
\hat{p} = a + b\, p,
\label{eq:prod}
\end{equation}
and the residual $y-\hat{p}$ was learned with the same covariates, learner, $\lambda$ and folds as in equation~(\ref{eq:residual}) (Supplementary Methods~S5).

\subsection{Yearly mapping, area of applicability and prediction intervals}
For Alaska and the Lena Delta, the recalibrated Stefan and physics-residual models were fitted to all cell-years of the region and applied to the 1 km cells with an ESA CCI permafrost probability of at least 50\% (\nncellmapAk{} and \nncellmapLe{} cells), using the full-season climate of each year from 2011 to 2024. The area of applicability followed Meyer and Pebesma\cite{MeyerPebesma2021}. With $\mathbf{z}$ the standardised 23-dimensional feature vector (with the 2011--2024 mean climate), the dissimilarity of map cell $k$ is $d_k = \min_{j} \lVert \mathbf{z}_k - \mathbf{z}_j \rVert_2$ over all training rows $j$, and the cell is inside the area of applicability when
\begin{equation}
d_k \leq \tau, \qquad \tau = Q_{0.95}\!\left\{ \min_{j:\, b(j) \neq b(i)} \lVert \mathbf{z}_i - \mathbf{z}_j \rVert_2 \;:\; i \in \text{training rows} \right\},
\label{eq:aoa}
\end{equation}
where $b(\cdot)$ denotes the spatial block. The 90\% prediction interval of each map value was taken from the 5th and 95th percentiles of the spatial-block out-of-fold residuals. For display, the 1 km values were interpolated bilinearly onto a 250 m hillshade grid.

\subsection{Computation}
Models were fitted with fixed random seeds on cloud nodes (Rescale; 96-core CPU nodes, and one NVIDIA A100 GPU node for the neural learners); data preparation, aggregation and figures were run on a local server (Supplementary Table~S4).

\subsection{Code availability}
The code for data preparation, modelling, evaluation and figures is available at [repository and DOI]. Every number in the text is generated from the stored result tables by the scripts that produce the manuscript macros.
